% ============================================================================
\section{Solving the general two-level problem}
\label{sec:geometry}
% ============================================================================

\subsection{Eigenvalues without diagonalizing}

The decomposition of \cref{sec:formalism} pays for itself immediately. Using the
product rule for Pauli matrices one finds (\cref{app:pauli})
%
\begin{equation}
\left(\dvec\cdot\pauliv\right)^2 = |\dvec|^2\,\Iop .
\label{eq:dsquared}
\end{equation}
%
A matrix whose square is a multiple of the identity can only have eigenvalues
$\pm|\dvec|$, so the spectrum of $\Hop$ follows with no linear algebra at all:

\begin{keyresult}[Splitting]
\vspace{-0.35cm}
\begin{equation}
E_{\pm} = \varepsilon_0 \pm \frac{|\dvec|}{2}
        = \varepsilon_0 \pm \half\sqrt{d_x^2+d_y^2+d_z^2},
\qquad
E_+ - E_- = |\dvec| .
\label{eq:eigenvalues}
\end{equation}
\end{keyresult}

\noindent
This one formula is the bonding--antibonding splitting, the generalized Rabi frequency,
the adiabatic gap at an avoided crossing, the Davydov splitting of an excitonic dimer,
and the inversion doublet of ammonia. It is worth pausing on what it says. Writing
$\dperp \equiv \sqrt{d_x^2+d_y^2}$ for the coupling magnitude,
%
\begin{equation}
E_+ - E_- = \sqrt{d_z^2 + \dperp^2} \;\geq\; \dperp .
\label{eq:noncrossing}
\end{equation}
%
The asymmetry and the coupling add \emph{in quadrature}, so the splitting can never fall
below the coupling. Even when the two basis states are exactly degenerate, they are
split by $\dperp$. This is the \textbf{non-crossing rule}: two states that are coupled
cannot become degenerate by tuning one parameter.

\subsection{The mixing angle}

Write $\dvec$ in spherical polar coordinates,
%
\begin{equation}
d_z = |\dvec|\cos\theta,\qquad
d_x = |\dvec|\sin\theta\cos\varphi,\qquad
d_y = |\dvec|\sin\theta\sin\varphi,
\end{equation}
%
so that
%
\begin{equation}
\tan\theta = \frac{\dperp}{d_z},
\qquad
\varphi = \atantwo(d_y,d_x).
\label{eq:mixing-angle}
\end{equation}
%
Substituting into \cref{eq:pauli-components} and solving (\cref{app:eigvec}) gives the
eigenvectors:

\begin{keyresult}[Eigenvectors]
\vspace{-0.35cm}
\begin{align}
\ket{+} &= \cos\tfrac{\theta}{2}\,\ket{\alpha} + e^{\ii\varphi}\sin\tfrac{\theta}{2}\,\ket{\beta},\\
\ket{-} &= -e^{-\ii\varphi}\sin\tfrac{\theta}{2}\,\ket{\alpha} + \cos\tfrac{\theta}{2}\,\ket{\beta}.
\label{eq:eigenvectors}
\end{align}
\end{keyresult}

\noindent
The angle $\theta$ is the single most useful number in two-level physics: it measures how
thoroughly the coupling has scrambled the basis states. The population of $\ket{\alpha}$
in the upper state is $\cos^2(\theta/2)$ and in the lower state $\sin^2(\theta/2)$.

\begin{center}
\small
\begin{tabular}{@{}lll@{}}
\toprule
Regime & $\theta$ & Eigenstates are\ldots \\
\midrule
$d_z \gg \dperp > 0$   & $\to 0$     & essentially $\ket{\alpha},\ket{\beta}$; coupling is a perturbation \\
$d_z = 0$              & $=\pi/2$    & exactly $50/50$; maximally delocalized \\
$-d_z \gg \dperp > 0$  & $\to \pi$   & localized again, with the labels exchanged \\
\bottomrule
\end{tabular}
\end{center}

\noindent
Note the half-angles in \cref{eq:eigenvectors}. As $d_z$ sweeps from $+\infty$ to
$-\infty$ at fixed $\dperp$, $\theta$ runs from $0$ to $\pi$ and $\ket{+}$ evolves
continuously from $\ket{\alpha}$ to $\ket{\beta}$. That continuous exchange of character,
without the levels ever touching, \emph{is} an avoided crossing. \Cref{fig:avoided}
shows both halves of the story.

% ---------------------------------------------------------------------------
\begin{figure}[t]
\centering
\begin{tikzpicture}
\begin{groupplot}[
  group style={group size=2 by 1, horizontal sep=1.6cm},
  notestyle, height=6.2cm,
  legend style={at={(0.5,-0.30)}, anchor=north, legend columns=2,
                draw=none, fill=none, font=\footnotesize},
  legend cell align=left,
]
% ---- left: the avoided crossing
\nextgroupplot[
  xlabel={$d_z$}, ylabel={energy},
  title={Levels repel},
  xmin=-6, xmax=6, ymin=-3.4, ymax=3.4,
]
\addplot[cMuted, dashed, thin, domain=-6:6, samples=2]{0.5*x};
\addlegendentry{uncoupled $H_{\alpha\alpha}$}
\addplot[cMuted, dotted, thin, domain=-6:6, samples=2]{-0.5*x};
\addlegendentry{uncoupled $H_{\beta\beta}$}
\addplot[cUpper, very thick, domain=-6:6, samples=200]{0.5*sqrt(x^2+1)};
\addlegendentry{$E_+$}
\addplot[cLower, very thick, domain=-6:6, samples=200]{-0.5*sqrt(x^2+1)};
\addlegendentry{$E_-$}
\draw[<->, cAccent, thick] (axis cs:0,-0.5) -- (axis cs:0,0.5);
\node[cAccent, anchor=west, font=\footnotesize] at (axis cs:0.35,0){$\dperp$};

% ---- right: the composition
\nextgroupplot[
  xlabel={$d_z$}, ylabel={population in $\ket{-}$},
  title={Character is exchanged},
  xmin=-6, xmax=6, ymin=-0.04, ymax=1.04,
]
\addplot[cLower, very thick, domain=-6:6, samples=300]
  {sin(0.5*atan2(1,x))^2};
\addlegendentry{$|\!\braket{\alpha|-}\!|^2$}
\addplot[cUpper, very thick, domain=-6:6, samples=300]
  {cos(0.5*atan2(1,x))^2};
\addlegendentry{$|\!\braket{\beta|-}\!|^2$}
\addplot[cMuted, dotted, thin, domain=-6:6, samples=2]{0.5};
\end{groupplot}
\end{tikzpicture}
\caption{The two faces of an avoided crossing, for fixed coupling $\dperp = 1$.
\textbf{Left:} the uncoupled diagonal energies (grey) cross at $d_z=0$; the
eigenvalues of \cref{eq:eigenvalues} do not, and are separated there by exactly
$\dperp$. \textbf{Right:} the composition of the lower eigenstate, from
\cref{eq:eigenvectors}. The exchange of character happens over a window of width
$\sim\dperp$ in $d_z$ --- a stronger coupling makes the switch \emph{more} gradual,
not less.}
\label{fig:avoided}
\end{figure}
% ---------------------------------------------------------------------------

\subsection{Dynamics: the Bloch equation}
\label{sec:bloch}

For a normalized state $\ket{\psi}$ define the \textbf{Bloch vector}
%
\begin{equation}
\svec = \Braket{\pauliv}
      = \Big(\Braket{\psi|\sxo|\psi},\;\Braket{\psi|\syo|\psi},\;\Braket{\psi|\szo|\psi}\Big).
\label{eq:bloch-vector}
\end{equation}
%
Three real numbers, and for any pure state $|\svec| = 1$ exactly, so the state lives on
the surface of a unit sphere. The north pole is $\ket{\alpha}$, the south pole is
$\ket{\beta}$, the equator holds the equally weighted superpositions, and longitude on
the equator is the relative phase. Comparing \cref{eq:bloch-vector} with
\cref{eq:eigenvectors}, the eigenstate $\ket{+}$ sits at polar angles $(\theta,\varphi)$
--- that is, $\svec_+ = \dhat$. \textbf{The eigenstates are the two points where the
state vector is parallel or antiparallel to $\dvec$.}

Applying the Heisenberg equation to each Pauli matrix and using
$[\hat\sigma_a,\hat\sigma_b]=2\ii\,\epsilon_{abc}\hat\sigma_c$ gives a strikingly simple
result, worked through in full in \cref{app:blocheq}:

\begin{keyresult}[Bloch equation]
\vspace{-0.35cm}
\begin{equation}
\frac{\dd\svec}{\dd t} = \frac{1}{\hbar}\,\dvec\times\svec .
\label{eq:bloch-eq}
\end{equation}
\end{keyresult}

\noindent
The Bloch vector \emph{precesses about $\dvec$} at angular frequency $|\dvec|/\hbar$.
Three consequences deserve to be stated explicitly.

\begin{enumerate}[leftmargin=*,itemsep=2pt]
\item \textbf{Energy is conserved.} A cross product is perpendicular to both arguments,
so $\dvec\cdot\svec$ is constant --- which is exactly the statement that
$\Braket{\Hop} = \varepsilon_0 + \half\dvec\cdot\svec$ does not change.

\item \textbf{Eigenstates do not move.} If $\svec \parallel \dvec$ the cross product
vanishes. Stationary states are stationary.

\item \textbf{The precession frequency is the Bohr frequency.} Since
$|\dvec| = E_+-E_-$, the rate is $(E_+-E_-)/\hbar$.
\end{enumerate}

\noindent
Notice that $\varepsilon_0$ is absent from \cref{eq:bloch-eq}. An overall energy shift
produces an overall phase, which no measurement on this system can detect.

% ---------------------------------------------------------------------------
\begin{figure}[t]
\centering
\tdplotsetmaincoords{72}{118}
\begin{subfigure}[b]{0.47\textwidth}
\centering
\begin{tikzpicture}[scale=2.05]
  % --- screen-space silhouette (NOT in 3d coords, so it stays a circle) ---
  \shade[ball color=gray!22, opacity=0.16] (0,0) circle (1);
  \draw[gray!50, thin] (0,0) circle (1);
  \begin{scope}[tdplot_main_coords]
    % equator and the xz meridian
    \tdplotdrawarc[gray!45,thin,densely dashed]{(0,0,0)}{1}{0}{360}{}{}
    \tdplotsetrotatedcoords{0}{90}{0}
    \tdplotdrawarc[tdplot_rotated_coords,gray!35,thin,densely dotted]{(0,0,0)}{1}{0}{360}{}{}
    \tdplotsetrotatedcoords{0}{0}{0}
    % axes
    \draw[-{Stealth[length=4pt]}, gray!65] (0,0,0) -- (1.38,0,0)
      node[anchor=north,black!65,font=\scriptsize]{$x$};
    \draw[-{Stealth[length=4pt]}, gray!65] (0,0,0) -- (0,1.38,0)
      node[anchor=west,black!65,font=\scriptsize]{$y$};
    \draw[-{Stealth[length=4pt]}, gray!65] (0,0,0) -- (0,0,1.30)
      node[anchor=south,black!65,font=\scriptsize]{$z$};
    % the d vector at theta = 54, phi = 40
    \tdplotsetcoord{P}{1}{54}{40}
    \tdplotsetcoord{Pxy}{0.809}{90}{40}
    \draw[gray!65, dashed, thin] (P) -- (Pxy);
    \draw[gray!65, dashed, thin] (0,0,0) -- (Pxy);
    \draw[-{Stealth[length=6pt]}, cAccent, very thick] (0,0,0) -- (P);
    \node[cAccent, anchor=south west, font=\small, inner sep=1pt] at (P) {$\dhat$};
    % theta arc, in the plane containing z and d
    \tdplotsetthetaplanecoords{40}
    \tdplotdrawarc[tdplot_rotated_coords,cAccent,thin,-{Stealth[length=3.5pt]}]%
      {(0,0,0)}{0.30}{0}{54}{anchor=east,font=\scriptsize,cAccent,inner sep=1.5pt}{$\theta$}
    % phi arc, in the xy plane
    \tdplotdrawarc[cUpper,thin,-{Stealth[length=3.5pt]}]{(0,0,0)}{0.52}{0}{40}%
      {anchor=north west,font=\scriptsize,cUpper,inner sep=1pt}{$\varphi$}
    % poles
    \fill[black] (0,0,1) circle (0.9pt);
    \node[anchor=east, font=\scriptsize, inner sep=2pt] at (0,0,1) {$\ket{\alpha}$};
    \fill[black] (0,0,-1) circle (0.9pt);
    \node[anchor=east, font=\scriptsize, inner sep=2pt] at (0,0,-1) {$\ket{\beta}$};
  \end{scope}
\end{tikzpicture}
\caption{Geometry. $\dvec$ has polar angles $(\theta,\varphi)$; the eigenstate
$\ket{+}$ points along $\dhat$ and $\ket{-}$ along $-\dhat$.}
\label{fig:bloch-geom}
\end{subfigure}
\hfill
\begin{subfigure}[b]{0.47\textwidth}
\centering
\begin{tikzpicture}[scale=2.05]
  \shade[ball color=gray!22, opacity=0.16] (0,0) circle (1);
  \draw[gray!50, thin] (0,0) circle (1);
  \begin{scope}[tdplot_main_coords]
    \tdplotdrawarc[gray!45,thin,densely dashed]{(0,0,0)}{1}{0}{360}{}{}
    \draw[-{Stealth[length=4pt]}, gray!65] (0,0,0) -- (1.38,0,0)
      node[anchor=north,black!65,font=\scriptsize]{$x$};
    \draw[-{Stealth[length=4pt]}, gray!65] (0,0,0) -- (0,1.38,0)
      node[anchor=west,black!65,font=\scriptsize]{$y$};
    \draw[-{Stealth[length=4pt]}, gray!65] (0,0,0) -- (0,0,1.30)
      node[anchor=south,black!65,font=\scriptsize]{$z$};
    % precession cone: half-angle 54.5 deg about d-hat, which lies in the xz plane
    \tdplotsetrotatedcoords{0}{54.5}{0}
    \tdplotdrawarc[tdplot_rotated_coords, cUpper, very thick]%
       {(0,0,0.5807)}{0.8141}{0}{360}{}{}
    \tdplotsetrotatedcoords{0}{0}{0}
    % the precession axis
    \tdplotsetcoord{D}{1.16}{54.5}{0}
    \draw[-{Stealth[length=6pt]}, cAccent, very thick] (0,0,0) -- (D);
    \node[cAccent, anchor=south east, font=\small, inner sep=3pt] at (D) {$\dhat$};
    % start of the trajectory, and a quarter period later
    % (s(pi/2) = s_par + d-hat x s_perp, worked out for a 54.5 deg cone)
    \fill[black] (0,0,1) circle (1.3pt);
    \fill[cUpper] (0.4728,-0.8141,0.3372) circle (1.3pt);
    \node[anchor=south, font=\scriptsize, inner sep=2pt] at (0,0,1.02) {$\svec(0)=\ket{\alpha}$};
    \fill[black] (0,0,-1) circle (0.9pt);
    \node[anchor=east, font=\scriptsize, inner sep=2pt] at (0,0,-1) {$\ket{\beta}$};
  \end{scope}
\end{tikzpicture}
\caption{Dynamics. Starting at $\ket{\alpha}$, the state sweeps a cone about $\dhat$
at rate $|\dvec|/\hbar$ (red dot: a quarter period later).}
\label{fig:bloch-prec}
\end{subfigure}
\caption{The Bloch sphere. Because the trajectory is a cone about $\dhat$, a state
prepared at the north pole reaches the south pole only if $\dhat$ lies in the
equatorial plane --- which is precisely the resonance condition of \cref{sec:driven}.}
\label{fig:bloch}
\end{figure}
% ---------------------------------------------------------------------------

\begin{nbbox}[In the notebook]
\Cref{eq:bloch-vector} is \texttt{bloch\_vector(psi)}; time evolution is
\texttt{propagate(H, psi0, times)}, which phases the eigenbasis coefficients rather
than integrating an ODE. The notebook checks \cref{eq:bloch-eq} indirectly, by
confirming that the projection $\svec\cdot\dhat$ is constant to machine precision
along a trajectory.
\end{nbbox}
